%! The Inverse and Determinant of a Matrix — Unit 4, Lesson 4.11 — Homework Key
\documentclass[10pt]{article}
\usepackage{precalculus-article}
\usepackage{precalculus-key}
\newcommand{\TallMath}[1]{$\displaystyle #1\rule[-1.4em]{0pt}{3.2em}$}

\begin{document}

\pageheader{Unit 4, Lesson 4.11}{Homework Key}
\namedateperiod

\vspace{0.1in}

\begin{notesbox}{Practice}
\begin{enumerate}[leftmargin=1.5em, itemsep=6pt]

    \item Compute the determinant of each matrix.
    \begin{enumerate}[label=(\alph*), leftmargin=1.5em, itemsep=10pt]
        \item $\begin{bmatrix}4&1\\2&3\end{bmatrix}$ \hfill $\det = \ans{$4(3)-1(2)=10$}$

        \item $\begin{bmatrix}-2&5\\1&-3\end{bmatrix}$ \hfill $\det = \ans{$(-2)(-3)-(5)(1)=6-5=1$}$

        \item $\begin{bmatrix}6&2\\3&1\end{bmatrix}$ \hfill $\det = \ans{$6(1)-2(3)=0$}$
    \end{enumerate}
    \vspace{0.05in}

    \item For each matrix in problem 1, determine whether it is invertible or singular.
          Show your reasoning.

    \vspace{4pt}
    \ans{(a) $\det=10\ne 0$ $\Rightarrow$ \textbf{invertible}.\\[3pt]
    (b) $\det=1\ne 0$ $\Rightarrow$ \textbf{invertible}.\\[3pt]
    (c) $\det=0$ $\Rightarrow$ \textbf{singular} (not invertible).}

    \vspace{0.05in}

    \item Find the inverse of each invertible matrix from problem 1.
          Show all work using the formula $A^{-1} = \dfrac{1}{\det(A)}\begin{bmatrix}d&-b\\-c&a\end{bmatrix}$.

    \vspace{4pt}
    \ans{(a) $A=\begin{bmatrix}4&1\\2&3\end{bmatrix}$,\quad $\det(A)=10$.\\[3pt]
    $A^{-1}=\dfrac{1}{10}\begin{bmatrix}3&-1\\-2&4\end{bmatrix}=\begin{bmatrix}3/10&-1/10\\-1/5&2/5\end{bmatrix}$.\\[6pt]
    (b) $B=\begin{bmatrix}-2&5\\1&-3\end{bmatrix}$,\quad $\det(B)=1$.\\[3pt]
    $B^{-1}=\dfrac{1}{1}\begin{bmatrix}-3&-5\\-1&-2\end{bmatrix}=\begin{bmatrix}-3&-5\\-1&-2\end{bmatrix}$.}

    \vspace{0.05in}

    \item Verify your answer to 3(a) by computing $A \cdot A^{-1}$ and confirming the result
          equals $I_2$.

    \vspace{4pt}
    \ans{$A\cdot A^{-1}=\dfrac{1}{10}\begin{bmatrix}4&1\\2&3\end{bmatrix}\begin{bmatrix}3&-1\\-2&4\end{bmatrix}=\dfrac{1}{10}\begin{bmatrix}12-2&-4+4\\6-6&-2+12\end{bmatrix}=\dfrac{1}{10}\begin{bmatrix}10&0\\0&10\end{bmatrix}=\begin{bmatrix}1&0\\0&1\end{bmatrix}=I$. \checkmark}

    \vspace{0.05in}

    \item The columns of $M = \begin{bmatrix}4&2\\6&3\end{bmatrix}$ form two vectors.
    \begin{enumerate}[label=(\alph*), leftmargin=1.5em, itemsep=10pt]
        \item Compute $\det(M)$.

        \vspace{4pt}
        \ans{$\det(M)=4(3)-2(6)=12-12=0$}

        \item Interpret $|\det(M)|$ geometrically: what does it tell you about the column vectors?

        \ansline{$|\det(M)|=0$ means the two column vectors are parallel (collinear) and span no area ---}
        \ansline{the parallelogram they form is completely flat, with area 0.}
    \end{enumerate}

    \item \textbf{AP-Style Multiple Choice.} Which of the following matrices does NOT have an
          inverse?
    \begin{enumerate}[label=(\Alph*), leftmargin=2em, itemsep=4pt]
        \item $\begin{bmatrix}1&0\\0&1\end{bmatrix}$
        \item $\begin{bmatrix}3&1\\6&2\end{bmatrix}$
        \item $\begin{bmatrix}2&-1\\3&4\end{bmatrix}$
        \item $\begin{bmatrix}5&2\\1&1\end{bmatrix}$
    \end{enumerate}

    \vspace{4pt}
    Answer: \ans{(B)} \quad
    \ans{$\det\!\begin{bmatrix}3&1\\6&2\end{bmatrix}=3(2)-1(6)=6-6=0$, so it has no inverse.}

\end{enumerate}
\end{notesbox}

\vspace{0.1in}

\begin{extensionbox}
    Let $A = \begin{bmatrix}a&0\\0&d\end{bmatrix}$ be a diagonal matrix. Derive a formula
    for $A^{-1}$ (assuming $a \ne 0$ and $d \ne 0$) and explain what it means geometrically.

    \vspace{4pt}
    \ans{$\det(A)=ad$.\quad $A^{-1}=\dfrac{1}{ad}\begin{bmatrix}d&0\\0&a\end{bmatrix}=\begin{bmatrix}1/a&0\\0&1/d\end{bmatrix}$.\\[3pt]
    Geometrically: $A$ scales the $x$-axis by $a$ and the $y$-axis by $d$; $A^{-1}$ undoes those scalings by the reciprocals.}
\end{extensionbox}

\end{document}
